\section{Evaluation}
\label{sec:evaluation}

\subsection{Setup, Baselines, and Metrics}
\label{sec:eval-method}

The evaluation asks whether SCOUT can make reliable decisions across the discrepancy patterns that arise in continuous heterogeneous validation, without depending on any particular
engine pair. We therefore organize the evaluation around concentrated faults, diffuse discrepancies, side imbalance (drift), and increasing discrepancy scale, rather than around a specific
cross-engine deployment.

The evaluation proceeds in two stages.
Parameterized synthetic workloads
(\S\ref{sec:eval-spatial}--\S\ref{sec:eval-operating})
isolate individual mechanisms and map operating boundaries.
YCSB Zipfian workloads (\S\ref{sec:eval-ycsb}) introduce
segment-size skew and hotness imbalance to test robustness
under more realistic data-distribution conditions.
Synthetic workloads cover concentrated, diffuse, and mixed
heavy-tail patterns with explicit sweeps over drift
intensity; exact constructions are given in
Supplement~\ref{sec:appendix-workloads}.

\paragraph{Baselines}
We compare three methods under the single-round fixed-budget
constraint.
B1 (per-segment checksum) maintains one checksum per segment;
because its digest grows with~$P$, it violates the
fixed-budget constraint and serves only as a full-information
reference rather than a same-budget baseline.
B2 (CountSketch ranking) treats each segment as an item and
outputs a top-$K$ candidate set ranked by difference intensity.
B3 (hierarchical bucket tree) maps segments to leaf nodes of a
fixed-depth tree and drills down from root to leaves by selecting
high-score buckets at each level; when a leaf bucket is selected,
all segments it contains become candidates, producing step-wise
cost jumps of $P/b^d$ per bucket
(Supplement~Fig.~\ref{fig:b3_bucket_tree}). B2 and B3 represent the two most natural fixed-budget single-round alternatives in this setting: unstructured ranking and structured hierarchical aggregation.
All methods are aligned by digest memory in bytes at a primary
tier of ${\sim}512$\,KiB; each baseline is additionally tested
at ${\sim}2{\times}$ memory
(exact configurations in Supplement~Table~\ref{tab:c5-0}).

\paragraph{Metrics and operating-point selection}
The primary triage metric is the candidate fraction
$\mathtt{cand\_frac}{=}|\mathcal{C}|/P$.
For triage comparisons, we sweep each method's internal tuning
knob and report the best operating point that satisfies the
target heavy-segment recall for the experiment, where
$\mathtt{recall\_big}$ measures recall over the truly heavily
dirty segments in the workload instance. Where a candidate-budget
cap is imposed, infeasible points are marked as \textbf{N/A}.

Gating quality is evaluated through false-alarm rate (FAR) and false-rejection rate (FRR) trade-offs. When
the experiment includes abstention, we additionally report
performance on decisive (non-\texttt{unknown}) samples.
Sizing accuracy is reported as $\mathrm{Std(rel\_err)}$, the
standard deviation of the relative estimation error
$(\hat{D}-D)/D$ across repetitions.

YCSB validation scales to ${\sim}2{\times}10^9$ records across
$P{=}100{,}000$ segments with one-sided digest below 1\,MiB. Details are in
Supplement~\ref{sec:appendix-scale}.
